% ============================================================================
% Fichier  : partie6/P06_C18_indicateurs.tex
% Titre    : Surveillance, Mesure et Indicateurs
% Auteur   : MINKA MI NGUIDJOÏ Thierry Emmanuel
% ============================================================================
% Sources : ISO/IEC 27004:2016 ; ISO/IEC 27001:2022 (clause 9.1). Fictif : Eséka.
% ============================================================================

\chapter{Surveillance, Mesure et Indicateurs}
\label{chap:indicateurs}
\chaptermark{Mesure et indicateurs}

\epigraph{%
    « Ce qui ne se mesure pas ne se pilote pas~;\\
    ce qui ne se pilote pas se dégrade. »%
}{Principe du cours}

\niveauChapitre{Intermédiaire}{La partie V.}

\begin{boxobjectifs}
\begin{itemize}
    \item Comprendre l'exigence de surveillance et de mesure (clause~9.1).
    \item Distinguer indicateurs de performance et de risque.
    \item Construire un indicateur utile (mesurable, suivi, décidable).
    \item Bâtir un tableau de bord de sécurité pour la direction.
    \item Relier la mesure à la décision et à l'amélioration.
\end{itemize}
\end{boxobjectifs}

\bigskip

% ============================================================================
\section*{Question directrice : comment savoir si le SMSI fonctionne vraiment ?}
% ============================================================================

La mairie a déployé des mesures. Fonctionnent-elles~? Sans mesure, nul ne le
sait. Ce chapitre construit les indicateurs et le tableau de bord qui
répondent.

\medskip
\textbf{Plan du chapitre.}
\begin{enumerate}
    \item L'exigence de mesure (section~\ref{sec:c18-exigence}).
    \item Performance et risque (section~\ref{sec:c18-kpikri}).
    \item Un bon indicateur (section~\ref{sec:c18-bon}).
    \item Le tableau de bord (section~\ref{sec:c18-tableau}).
    \item De la mesure à la décision (section~\ref{sec:c18-decision}).
\end{enumerate}

\separateur

% ============================================================================
\section{L'Exigence de Mesure}
\label{sec:c18-exigence}
% ============================================================================

\begin{boxdef}
La clause~9.1 de \norme{ISO/IEC 27001} exige de déterminer \emph{quoi}
surveiller et mesurer, \emph{comment}, \emph{quand} et \emph{qui} analyse les
résultats. La norme \norme{ISO/IEC 27004} en donne les lignes directrices.
\end{boxdef}

Mesurer n'est pas accumuler des chiffres~: c'est choisir les quelques
indicateurs qui éclairent une décision. Un bon dispositif de mesure est
sélectif.

% ============================================================================
\section{Performance et Risque}
\label{sec:c18-kpikri}
% ============================================================================

\begin{boxdef}
Un \termecle{indicateur de performance} (\sigle{KPI}) mesure l'efficacité d'une
mesure (taux d'agents formés). Un \termecle{indicateur de risque} (\sigle{KRI})
alerte sur une exposition croissante (nombre de comptes sans \sigle{MFA}).
\end{boxdef}

Les deux se complètent~: les \sigle{KPI} disent « où en est-on », les
\sigle{KRI} disent « où le danger monte ». Un tableau de bord équilibré
comporte les deux.

% ============================================================================
\section{Un Bon Indicateur}
\label{sec:c18-bon}
% ============================================================================

\begin{tableaucours}{L{3.6cm}L{8.6cm}}{%
    \tableauHeader{Qualité} &
    \tableauHeader{Signification}
}
Mesurable & On peut le calculer objectivement \\
Pertinent & Il éclaire une décision réelle \\
Suivi & Il est relevé à intervalle régulier \\
Doté d'une cible & On sait ce qui est bon ou mauvais \\
Attribué & Quelqu'un en répond \\
\end{tableaucours}
\vspace{-0.8em}
\begin{center}{\small\itshape Tableau~18.1. Qualités d'un bon indicateur}\end{center}

Un indicateur que personne ne regarde ou qui ne change aucune décision est à
supprimer~: il coûte sans servir.

% ============================================================================
\section{Le Tableau de Bord}
\label{sec:c18-tableau}
% ============================================================================

\begin{tableaucours}{L{4.6cm}L{2.2cm}L{2.2cm}L{2.0cm}}{%
    \tableauHeader{Indicateur} &
    \tableauHeader{Type} &
    \tableauHeader{Cible} &
    \tableauHeader{État}
}
Agents formés & KPI & 100 \% & 80 \% \\
Sauvegardes testées & KPI & Mensuel & À jour \\
Comptes sans MFA & KRI & 0 & 5 \\
Incidents du mois & KRI & -- & 2 \\
\end{tableaucours}
\vspace{-0.8em}
\begin{center}{\small\itshape Tableau~18.2. Extrait du tableau de bord de la mairie}\end{center}

\begin{boxmetier}{RSSI}
\textbf{Un tableau de bord pour décider.} La RSSI présente au comité un tableau
d'une page, lisible par un non-technicien. Chaque indicateur hors cible
appelle une décision~: former les agents restants, activer la MFA manquante.
Le tableau de bord n'est pas un rapport~: c'est un outil de pilotage.
\end{boxmetier}

% ============================================================================
\section{De la Mesure à la Décision}
\label{sec:c18-decision}
% ============================================================================

\begin{boxprocedure}
\textbf{Constatation.} Le tableau de bord montre 5~comptes sans \sigle{MFA} et
un taux de formation de 80~\%.

\medskip
\textbf{Analyse.} Deux écarts à la cible appellent deux décisions~: activer la
MFA sur les cinq comptes, achever la formation des agents restants. La mesure
débouche sur l'action~: c'est le lien entre « Check » et « Act » du
\sigle{PDCA} (Chapitre~\ref{chap:pdca}).
\end{boxprocedure}

La mesure alimente la revue de direction et l'audit (Chapitre~\ref{chap:audit})~:
elle fournit les preuves chiffrées que le \sigle{SMSI} fonctionne et s'améliore.

\separateur

\begin{boxresume}
\textbf{Ce qu'il faut retenir de ce chapitre~:}
\begin{itemize}
    \item La \textbf{clause~9.1} exige de mesurer quoi, comment, quand et par
          qui.
    \item Les \textbf{KPI} mesurent la performance~; les \textbf{KRI} alertent
          sur le risque.
    \item Un \textbf{bon indicateur} est mesurable, pertinent, suivi, doté
          d'une cible, attribué.
    \item Le \textbf{tableau de bord} est un outil de pilotage, lisible par la
          direction.
    \item La \textbf{mesure débouche sur la décision} et nourrit l'audit.
\end{itemize}
\end{boxresume}

\begin{boxinfo}
\textbf{Points de contrôle de l'élève.} Je sais~: énoncer l'exigence de la
clause~9.1~; distinguer KPI et KRI~; reconnaître un bon indicateur~; bâtir un
tableau de bord~; relier mesure et décision.
\end{boxinfo}

\bigskip

\begin{boxexo}[title={\ding{45}\quad Exercice~18.1. Quoi mesurer \niveaubase}]
Que doit déterminer la clause~9.1 au sujet de la mesure~?
\end{boxexo}

\begin{boxexo}[title={\ding{45}\quad Exercice~18.2. KPI ou KRI \niveaubase}]
« Taux d'agents formés » et « comptes sans MFA »~: lequel est un KPI, lequel un
KRI~?
\end{boxexo}

\begin{boxexo}[title={\ding{45}\quad Exercice~18.3. Bon indicateur \niveaubase}]
Citez trois qualités d'un bon indicateur.
\end{boxexo}

\begin{boxexo}[title={\ding{45}\quad Exercice~18.4. Indicateur inutile \niveaumoyen}]
Pourquoi supprimer un indicateur que personne ne regarde~?
\end{boxexo}

\begin{boxexo}[title={\ding{45}\quad Exercice~18.5. Construire un KPI \niveaumoyen}]
Construisez un KPI pour la mesure « sauvegarde de l'état civil ».
\end{boxexo}

\begin{boxexo}[title={\ding{45}\quad Exercice~18.6. Hors cible \niveaumoyen}]
Pour chaque indicateur hors cible du tableau de bord, proposez une décision.
\end{boxexo}

\begin{boxexo}[title={\ding{45}\quad Exercice~18.7. Lisible \niveaumoyen}]
Pourquoi le tableau de bord doit-il être lisible par un non-technicien~?
\end{boxexo}

\begin{boxexo}[title={\ding{45}\quad Exercice~18.8. Check et Act \niveaumoyen}]
Montrez que la mesure relie les étapes Check et Act du PDCA.
\end{boxexo}

\begin{boxexo}[title={\ding{45}\quad Exercice~18.9. Tableau de bord \niveauexpert}]
Construisez un tableau de bord de six indicateurs pour la mairie d'Eséka (type,
cible, responsable).
\end{boxexo}

\begin{boxexo}[title={\ding{45}\quad Exercice~18.10. Note de pilotage \niveauexpert}]
Rédigez la note mensuelle de pilotage que la RSSI présente au comité de
sécurité.
\end{boxexo}

\bigskip

\begin{boxinfo}
\textbf{Pour aller plus loin.}
\begin{itemize}
    \item ISO/IEC 27004~\cite{iso27004}~; ISO/IEC 27001, clause~9.1~\cite{iso27001}.
\end{itemize}
\end{boxinfo}

\begin{boxinfo}
\textbf{Transition.} Mesurer soi-même ne suffit pas~: il faut un regard
extérieur. Le Chapitre~\ref{chap:audit} traite l'audit et la certification.
\end{boxinfo}
